\subsection{格与序}

令 V 为一个 Q-向量空间。V 中的格是 V 的一个自由 Z-子模 V，使得由 V 嵌入 V 所诱导的 Q-线性映射 \(\alpha_{\mathcal{V},V} : \mathcal{V} \otimes_{Z} Q \to V\) 为双射。当 V 有限维时，这等价于说 V 是有限型的 Z-子模并生成 Q-向量空间 V（回忆一下，有限型的无挠 Z-模是自由的，见《代数》第 VII 章第4节第4号定理的推论2）；此外，此时我们的定义是《交换代数》第 VII 章第4节第1号定义1的特例（同上，例3）。若 W 是 V 的向量子空间，且 V 是 V 中的一个格，则 \(V \cap W\) 是 W 中的格。

若 V 是 Q-代数，则 V 中的一个序是底层向量空间中的格 V，且是 V 的 Z-子代数；此时映射 \(\alpha_{V,V}\) 是 Q-代数的同构。若 V 是有单位的 Q-代数，则 V 中的幺序是包含单位元的 V 中的序。

假设 V 是一个 Q-双代数，余积为 c，余单位为 \(\gamma\)。若 V 是向量空间 V 中的一个格，则典则映射 \(i: V \otimes_{Z} V \to V \otimes_{Q} V\) 是单射；V 中的双序是有单位代数 V 中的一个幺序 V，使得 \(\gamma(\mathcal{V}) \subset \mathbf{Z}\) 且 \(c(\mathcal{V}) \subset i(\mathcal{V} \otimes_{\mathbf{Z}} \mathcal{V})\)；映射

\[
\gamma_ {\mathcal {V}} \colon \mathcal {V} \to \mathbf {Z} \text { and } c _ {\mathcal {V}} \colon \mathcal {V} \to \mathcal {V} \otimes_ {\mathbf {Z}} \mathcal {V}
\]

由 \(\gamma\) 和 c 诱导的  及  赋予 V 以 Z-双代数结构，而映射 \(\alpha_{\mathcal{V},V}\) 此时是 Q-双代数的同构。

